\documentclass[11pt]{article}
\usepackage[a4paper,margin=25mm]{geometry}
\usepackage{amsmath,amssymb,amsthm,hyperref}
\newtheorem{theorem}{Theorem}
\newtheorem{lemma}{Lemma}
\newtheorem{corollary}{Corollary}
\newcommand{\PP}{\mathcal P}
\title{Polynomial conditioning and rational ranks for two-variable ordered prefixes}
\author{Proof independently reviewed by size\_bounds}
\date{30 September 2026}
\begin{document}\maketitle
Write $w=t(1-t)$, $H_d(t)=\int_0^tP_d(2u-1)du$, and
$R_i(t)=t-\eta_i H_{d_i}(t)$, where $0<\eta_i\le1/2$.
The increasing functions $R_i$ are distribution functions. Define
\[
 J(t)=\int_0^t R_1(u)R_2'(u)\,du.
\]
Thus $J(t)$ is the probability that independent variables with these
CDFs satisfy $X_1<X_2<t$. Its reversed counterpart is $R_1R_2-J$.

\begin{lemma}\label{condition}
Let $r\ge0$ and choose
\[
 d_1>2(r+2)+1,\qquad d_2>2(d_1+r+3)+1,
 \qquad D=r+d_1+d_2+2.
\]
For $p_0,p_1,p_2,p_{12},p_J\in\PP_r$,
\[
 \left(\|p_0\|_2^2+\|p_1\|_2^2+\|p_2\|_2^2+
 \|p_{12}\|_2^2+\|p_J\|_2^2\right)^{1/2}
 \le C(D+1)^{20}(\eta_1\eta_2)^{-2}
 \|p_0+p_1R_1+p_2R_2+p_{12}R_1R_2+p_JJ\|_2.
\]
All norms use Lebesgue measure on $[0,1]$.
\end{lemma}
\begin{proof}
We use the exact primitive and trinary-band lemmas from
\path{ordered_prefix_legendre_tools.tex}.
Set $\lambda_1=d_1(d_1+1)$ and $I_1(t)=\int_0^tH_{d_1}(u)du$.
The primitive formula gives
\[
 I_1=\frac{w'H_{d_1}-wH_{d_1}'}{\lambda_1-2}.
\]
Expanding $J$ yields
\[
 J=t^2/2-\eta_1I_1
     -\eta_2\int_0^t uH_{d_2}'(u)du
     +\eta_1\eta_2\int_0^tH_{d_1}(u)H_{d_2}'(u)du.
\]
The last two terms have the form $A H_{d_2}+B H_{d_2}'$ with
$\deg A,\deg B\le d_1+2$. This follows from the primitive lemma
with coefficient polynomials $u$ and $H_{d_1}$, respectively.
The other terms containing $H_{d_2}$ in the stated linear combination
have coefficient degrees no larger. After multiplication by the $p$'s,
all such high-band terms are orthogonal to $\PP_{d_1+r+2}$ by the
choice of $d_2$.

Consequently the orthogonal projection $L$ of the full combination
$F$ onto that low-degree space is exactly
\[
 L=p_0+t(p_1+p_2)+t^2p_{12}+\tfrac12t^2p_J
       -\eta_1(p_1+tp_{12})H_{d_1}-\eta_1p_JI_1.
\]
In its trinary representation $P+H_{d_1}Q+H_{d_1}'V$, the coefficient
\[
 V=\frac{\eta_1}{\lambda_1-2}w p_J
\]
is unique, and all three coefficient degrees are at most $r+2$.
Trinary coercivity bounds $\|V\|_2\le C(D+1)^6\|L\|_2$.
For a polynomial $p\in\PP_r$,
\begin{equation}\label{wbound}
 \|wp\|_2\ge c(r+1)^{-2}\|p\|_2.
\end{equation}
Indeed, exclude endpoint intervals of length $c(r+1)^{-2}$; the
polynomial evaluation bound preserves at least half the squared norm,
and $w$ is at least $c(r+1)^{-2}$ on what remains.
It follows that
\[
 \|p_J\|_2\le C(D+1)^{10}\eta_1^{-1}\|F\|_2.
\]
Subtract $p_JJ$. Since $0\le J\le1$, the remaining combination
has norm at most $(1+C(D+1)^{10}\eta_1^{-1})\|F\|_2$.
The two-band rank-product conditioning theorem from
\path{lacunary_legendre_conditioning.tex} bounds its four coefficients
by $C(D+1)^{10}(\eta_1\eta_2)^{-1}$ times this norm. The displayed
bound follows, after increasing the absolute constant.
\end{proof}

\begin{theorem}[Rational ordered-prefix data]\label{ranks}
For every fixed $T\ge1$ and every $m\ge1$, there are
$K=O((m+1)^4)$ separated ordered real nodes $b_q\in(0,1)$ and rational
arrays $r_{1,q},r_{2,q},a_q$ with a common denominator bounded by a
fixed power of $m+1$ (depending on $T$), such that for $0\le j\le m$,
\begin{align*}
 K^{-1}\sum_q b_q^j&=\frac1{j+1},&
 K^{-1}\sum_q b_q^j r_{i,q}&=\frac1{j+2}\quad(i=1,2),\\
 K^{-1}\sum_q b_q^j r_{1,q}r_{2,q}&=\frac1{j+3},&
 K^{-1}\sum_q b_q^j a_q&=\frac1{2(j+3)}.
\end{align*}
The node and rank gaps, including their endpoint gaps, are at least
$c/K$, and $|r_{i,q}-b_q|\le C_T(m+1)^{-T}$.
Furthermore these rational arrays are locally consistent ordered-prefix
data: append
\[
 (r_{1,0},r_{2,0},a_0)=(0,0,0),\qquad
 (r_{1,K+1},r_{2,K+1},a_{K+1})=(1,1,1/2).
\]
For every $0\le q\le K$,
\begin{equation}\label{area}
 0<a_{q+1}-a_q-r_{1,q}(r_{2,q+1}-r_{2,q})
 <(r_{1,q+1}-r_{1,q})(r_{2,q+1}-r_{2,q}).
\end{equation}
This does not assert a full fair family with three small dice.
\end{theorem}
\begin{proof}
Choose $d_1,d_2=O(m+1)$ satisfying Lemma~\ref{condition} for $r=m-1$,
and, if needed, increase their absolute additive constants so that
$d_1\ge m+3$ and $d_2>m+d_1+3$.
Choose $\eta_1=\eta_2=c(m+1)^{-(T+1)}$.
Legendre orthogonality gives, for $0\le j\le m$,
\[
 \int t^jR_i(t)dt=\frac1{j+2},\quad
 \int t^jR_1(t)R_2(t)dt=\frac1{j+3},\quad
 \int t^jJ(t)dt=\frac1{2(j+3)}.
\]
For the last identity, integration by parts gives
$[J(1)-\int t^{j+1}R_1R_2']/({j+1})$; orthogonality gives
$J(1)=1/2$ and the other integral $1/(j+3)$.
Take a real equal-weight quadrature of sufficiently high degree
$O(m)$ with $K=O(m^4)$ near-midpoint nodes $b_q^0$, as in the reviewed
separated quadrature lemma. Its degree is at least twice the maximal
degree of a polynomial in Lemma~\ref{condition}, ensuring exact Gram
matrices as well as all the displayed moment identities.

Round and correct the two rank columns to polynomial-denominator
rationals, preserving exactly their three constant mixed moments, by
the disjoint-row construction in
\path{../../independent_directions/multirank_step_compatibility.tex}.
The correction from the real baseline can be made $O((m+1)^{-A})$
for any fixed $A$ by increasing the initial polynomial denominator.
Independently round the $a_q=J(b_q^0)$ to a polynomial grid whose
denominator is a multiple of six and of the rank denominator.
Flooring and adding one grid unit to a suitable number of entries
preserves their exact mean $1/6$, with at most one grid unit error
per entry.

Keep all rational arrays fixed, and correct the nodes to restore the
nonconstant moment equations. In an antiderivative Legendre basis,
the baseline Jacobian evaluates $\phi_j$ times
$1,R_1,R_2,R_1R_2,J$, $0\le j<m$. Its empirical Gram is exact, and
Lemma~\ref{condition} gives a polynomial least singular value.
All derivative bounds are polynomial. Thus the same fixed-inverse
Newton contraction used in the multirank theorem moves the nodes by
$O((m+1)^{-A'})$, with arbitrarily prescribed fixed $A'$, if the
initial denominator exponent is sufficiently large. This keeps all
gaps and gives the stated exact moments and closeness.

Finally, at the continuous baseline the middle expression in
\eqref{area} equals
\[
 \int_{b_q^0}^{b_{q+1}^0}
       [R_1(t)-R_1(b_q^0)]R_2'(t)\,dt.
\]
Both this and its complement to the product of the two rank increments
are at least $c(b_{q+1}^0-b_q^0)^2\ge c/K^2$, since
$1/2\le R_i'\le3/2$. Include the endpoint intervals in this argument.
Errors much smaller than $K^{-2}$ preserve the strict inequalities.
Moving the nodes does not change them. All rank arrays retain one
polynomial common denominator.
\end{proof}

\paragraph{A related first-moment variant.}
The same argument conditions $1,R,M$, where
$R=t-\eta H_d$ and
\[
 M(t)=\int_0^t uR'(u)du
 =t^2/2-\eta\left[tH_d-
       \frac{w'H_d-wH_d'}{d(d+1)-2}\right].
\]
The unique $H_d'$ coefficient again isolates the coefficient of $M$
via~\eqref{wbound}. This can rationalize both cumulative mass and
cumulative first moment while preserving their mixed scalar moments.
It does not force the corrected real nodes to have polynomial rational
denominators, nor does it give positive completion with additional
independent variables.
\end{document}
